\scanpage{48}
\mspage{39}

$2'')$ $\forall i\in I$, posant $I'(i)=\{j\in I'\mid X_j^{\circ}\subset X_i\}$\unsure{-,} \deleted{on a}
\[
\deleted{I(i)=\{j\in I\mid X_j^{\circ}\subset X_i\}}
\]
\deleted{$\unsure{X_j}$}
\deleted{(de sorte que $I'(i)=I(i)\cap I'$, et $\bigcup_{j\in I(i)}X_j^{\circ}=\bigcup_{j\in I'(\unsure{j|i})}X_j^{\circ}$)}
on a \deleted{\unsure{que}}
\[
X_i=\bigcup_{j\in I'(i)}X_j^{\circ}
\]
\deleted{ou ce qui revient au même}
\[
\deleted{X_i=\bigcup_{j\in I(i)}X_j^{\circ}.\quad\text{i.e. }3')}
\]

$3'')$ On a, pour tout $i\in I$
\[
I'(i)=I'\deleted{\unsure{(i)}}\cap I_{\leq i}
\]
\margin{ce qui implique $\bigcup_{j\in I'(i)}X_j^{\circ}=\bigcup_{j\leq i}X_j^{\circ}$}
On voit donc que compte tenu de $3''$, $2'')$ équivaut à :

Ceci dit, l'ens. des conditions $1''$, $2''$, $3''$ équivaut aussi à \deleted{$1'',2'',3''$,} celui-ci
\[
\left\{\begin{aligned}
&\deleted{\unsure{\forall i\in I}}\quad|\Phi|=\bigcup_{i\in I}X_i^{\circ}&&\text{i.e. }(1')\\
&\forall i,j\in I\quad X_i^{\circ}\cap X_j^{\circ}=\varnothing&&\text{i.e. }(2')\\
&\forall i\in I\quad X_i=\bigcup_{j\leq i}X_j^{\circ}&&\text{i.e. }(3')\\
&+\ \text{la condition supplémentaire }3'')\Longleftrightarrow\gamma)\\
&\forall i,j\in I\quad X_j^{\circ}\subset X_i\ \inserted{\text{et }X_j^{\circ}\ne\varnothing}\Longrightarrow j\leq i.
\end{aligned}\right.
\]
\margin{NB}
Cela montre donc que
\[
(iv)\Longleftrightarrow(i)+\gamma)
\]
et l'équivalence $(iv)\Longleftrightarrow(i)$ revient à $(i)\Rightarrow(\gamma)$ qui est contenue dans le

\uline{Corollaire 1} Sous les conditions équivalentes (i), (iii), on a, pour $i,j\in I$, \deleted{\unsure{les}} \inserted{\uuline{tel que} $X_j^{\circ}\ne0$} \unsure{les} \deleted{\unsure{équivalences entre les}} \deleted{conditions suivantes}
\[
\deleted{\unsure{\forall}}(j\leq\deleted{\unsure{j}}i)\Longleftrightarrow(X_j\subset X_{\deleted{\unsure{j}}i})\Longleftrightarrow(X_j^{\circ}\subset X_{\deleted{\unsure{j}}i})\overset{\deleted{\unsure{\forall i,j}}\ \unsure{\text{i.e. }X_j^{\circ}\ne\varnothing}}{\Longleftrightarrow}(X_j^{\circ}\cap \unsure{X_i|X_i^{\circ}}\ne\varnothing)
\]
(\uline{NB} Dans l'énoncé de la prop, on n'a pas supposé que $i\leq j\Rightarrow X_i\subset X_j$ \unsure{??!})
% margin note written slanted
\margin{Attention en général (si on ne suppose $X_j^{\circ}\ne\varnothing$) la condition $X_j\subset X_i$ \unsure{est impliquée} par $j\leq i$. (Prendre tous les $X_i$ vides ! La condition (i) est \unsure{tautologiquement|tautologique et} vérifiée...)}

\uline{Dém du Cor.~1} en vertu de $3')$ on a
\[
X_i=\bigcup_{i'\leq i}X_{i'}^{\circ},\quad X_j=\bigcup_{j'\leq j}X_{j'}^{\circ}
\]
donc
\[
i\leq j\Longrightarrow X_i\subset X_j\quad\unsure{\longleftarrow}\ (\text{sans supposer }X_j^{\circ}\ne\varnothing)
\]
% archive sheet number: 48
